% =============================================================================
% introduction.tex — Project-level introduction chapter
% This chapter is NOT a module. It describes the system as a whole.
% =============================================================================

\chapter*{Project Introduction}
\addcontentsline{toc}{chapter}{Project Introduction}
\markboth{Project Introduction}{}

% -----------------------------------------------------------------------------
\section*{What Craftalism Is}
% -----------------------------------------------------------------------------

Craftalism is a modular, production-oriented economy platform that integrates a Minecraft game server with a fully externalized web backend. Rather than storing economy state inside the game server — the conventional approach for Minecraft plugins — the platform delegates all balance and transaction data to a centralized REST API. The Minecraft plugin becomes a thin, authenticated client: it handles commands and player interaction, then delegates all state reads and writes to the backend over HTTP.

This architecture allows the economy to be administered, audited, and observed entirely outside the game, through a web dashboard, while the game server itself remains stateless with respect to financial data.

% -----------------------------------------------------------------------------
\section*{System Architecture}
% -----------------------------------------------------------------------------

The platform is composed of five independent services, each in its own repository. All inter-service communication uses HTTP with OAuth2 bearer tokens.

\begin{figure}[h]
\centering
\begin{tikzpicture}[node distance=1.6cm and 2.4cm]

  % Nodes
  \node[external]                         (minecraft) {Minecraft\\Server\\(Plugin)};
  \node[service, right=of minecraft]      (api)       {Craftalism\\API};
  \node[database, right=of api]           (db)        {PostgreSQL};
  \node[service, below=of api]            (auth)      {Authorization\\Server};
  \node[external, right=of auth]          (jwks)      {JWKS\\(public keys)};
  \node[service, above=of api, yshift=0.4cm] (dashboard) {Dashboard};

  % Edges
  \draw[flow] (minecraft) -- node[flowlabel, above]{REST + JWT} (api);
  \draw[flow] (api)       -- node[flowlabel, above]{reads/writes} (db);
  \draw[flow] (auth)      -- (db);
  \draw[flow] (minecraft) -- node[flowlabel, left]{token request} (auth);
  \draw[flow] (api)       -- node[flowlabel, right]{validate via} (jwks);
  \draw[flow] (auth)      -- node[flowlabel, below]{publishes} (jwks);
  \draw[flow] (dashboard) -- node[flowlabel, right]{REST} (api);

\end{tikzpicture}
\caption{Craftalism system topology}
\label{fig:system-topology}
\end{figure}

\begin{enumerate}
  \item The Minecraft plugin authenticates with the \textbf{Authorization Server} on startup using the \code{client\_credentials} OAuth2 grant, and caches the resulting JWT access token.
  \item All subsequent plugin calls to the \textbf{Craftalism API} carry that JWT as a \code{Bearer} token in the \code{Authorization} header.
  \item The API validates tokens \emph{locally} by fetching the Authorization Server's public keys from \code{/oauth2/jwks} — no per-request round-trip to the auth server is required.
  \item The \textbf{Dashboard} calls the same API through an Nginx reverse proxy and renders the data for administrators.
  \item Both the API and the Authorization Server persist state to a shared \textbf{PostgreSQL} instance, in separate databases (\code{craftalism} and \code{authserver}).
\end{enumerate}

% -----------------------------------------------------------------------------
\section*{Repository Index}
% -----------------------------------------------------------------------------

\rowcolors{2}{CraftLight}{white}
\begin{tabularx}{\linewidth}{@{} L{3.4cm} L{5.0cm} X @{}}
  \toprule
  {\tableheadstyle Repository} &
  {\tableheadstyle Primary Stack} &
  {\tableheadstyle Purpose} \\
  \midrule
  \code{craftalism-api}                  & Java 17, Spring Boot 3.5, JPA   & Core REST API for players, balances, and transactions. \\
  \code{craftalism-authorization-server} & Java 17, Spring Authorization Server & OAuth2/OIDC token issuer for internal services. \\
  \code{craftalism-economy}              & Java 21, Paper API 1.21.4       & Minecraft plugin; delegates all economy state to the API. \\
  \code{craftalism-dashboard}            & React 19, TypeScript 5, Vite 7  & Admin web UI for read-oriented data visibility. \\
  \code{craftalism-deployment}           & Docker Compose, PostgreSQL 18   & Container orchestration for the full platform. \\
  \bottomrule
\end{tabularx}

% -----------------------------------------------------------------------------
\section*{Shared Conventions}
% -----------------------------------------------------------------------------

The following conventions apply across all modules and are referenced throughout this document.

\begin{itemize}
  \item \textbf{Authentication:} All API traffic from internal services is authenticated with a short-lived JWT issued via the \code{client\_credentials} grant. Tokens are validated locally using the published JWKS.
  \item \textbf{Error responses:} The API returns RFC 9457 \code{ProblemDetail} payloads for all error conditions, with a consistent set of extension fields: \code{type}, \code{detail}, \code{status}, \code{timestamp}, \code{path}, and \code{errors}.
  \item \textbf{Scope enforcement:} Read operations require the \code{api:read} scope. Write operations require \code{api:write}.
  \item \textbf{Containerization:} All services are distributed as Docker images and orchestrated with Docker Compose via the \code{craftalism-deployment} repository.
\end{itemize}
